%% Elsevier CAS single-column template
%% Manuscript: Adjoint transport theory for rare-event sampling in spin systems
%% Structure mirrors brks_article_.tex exactly:
%%   preamble -> frontmatter (shorttitle/shortauthors/title/authors/affiliations/
%%   credit/abstract/highlights/keywords/maketitle) -> Introduction ->
%%   Theoretical Approach -> Results and Discussion -> Conclusions ->
%%   Acknowledgements -> \printcredits -> bibliography.
%% Floats follow the brks convention: figures as non-floating minipage+\captionof,
%% tables as [H], \FloatBarrier before each new top-level section.

\documentclass[a4paper,fleqn]{cas-sc}

%% ---- Citation / bibliography style -------------------------------------
\usepackage[numbers,sort&compress]{natbib}
\usepackage{rotating}
\usepackage{booktabs}
\usepackage{amsmath}
\usepackage{graphicx}
\usepackage{caption}   % enables \captionof for non-floating figures
\usepackage{float}     % enables [H] placement option
\usepackage{placeins}  % \FloatBarrier confines floats to their section

%% ---- Packages carried over from the original manuscript ----------------
\usepackage{amssymb}
\usepackage{enumitem}
\usepackage{xcolor}
\usepackage{tikz}
\usetikzlibrary{arrows.meta,decorations.pathreplacing}
\definecolor{acc}{HTML}{2c5aa0}
\setlist{itemsep=2pt,topsep=3pt}
\newtheorem{claim}{Claim}
\newtheorem{prop}{Proposition}
\newcommand{\dg}{^{\dagger}}
\newcommand{\ip}[2]{\langle #1,\,#2\rangle}
\newcommand{\E}{\mathbb{E}}

%% ---- Confine floats to the subsection that introduces them -------------
%% Ensures every [H] table stays under its own subsection heading.
\let\Oldsubsection\subsection
\renewcommand{\subsection}{\FloatBarrier\Oldsubsection}

\begin{document}
	\let\WriteBookmarks\relax
	\def\floatpagepagefraction{1}
	\def\textpagefraction{.001}
	
	% Short title
	\shorttitle{Adjoint importance for rare-event sampling in spin systems}
	
	% Short author
	\shortauthors{M.Y. Alzaatreh}
	
	% Main title
	\title[mode=title]{From Radiation Shielding to Spin Glasses: Auditing the Adjoint-Importance Toolkit for Rare-Event Sampling}
	
	% ---- Authors -----------------------------------------------------------
	\author{M.Y. Alzaatreh}
	\cormark[1]
	\ead{mzaatreh@fbsu.edu.sa}
	\credit{Conceptualization, Methodology, Software, Formal analysis, Investigation, Writing -- original draft, Writing -- review \& editing}
	
	\affiliation{organization={Natural Sciences Unit, Fahad Bin Sultan University},
		city={Tabuk},
		country={Saudi Arabia}}
	
	\cortext[1]{Corresponding author}
	
	
	% ---- Abstract ----------------------------------------------------------
	\begin{abstract}
		Radiation transport's adjoint importance function and CADIS/weight window machinery mathematically identical to the rare-event committor are transferred to statistical mechanics and tested end-to-end on the 2D Ising ferromagnet and Edwards–Anderson spin glass. The adjoint committor identity is verified to machine precision, and Weighted Ensemble sampling (weight windows on spin configurations) reaches two decades deeper into the tail than naive Monte Carlo at matched cost. A self consistency training objective shaped directly on the harmonic committor condition finds rare events significantly more often than a hand-picked energy coordinate (66.7\% vs. 40.0\% hit rate, p = 0.017). A global variance reduction objective (FW-CADIS) achieves an eightfold cost reduction while matching uniform allocation's error profile. These results demonstrate the transport toolkit's effectiveness in the genuinely rare regime, and highlight the importance of a committor-shaped training objective rather than a convenient proxy.
	\end{abstract}
	% ---- Highlights ---------------------------------------------------------
	\begin{highlights}
		\item Transport adjoint importance and the rare-event committor are identical; verified to machine precision.
		\item Weighted Ensemble validated against exact enumeration; two decades deeper than naive Monte Carlo.
		\item FW-CADIS attains equivalent threshold coverage at roughly one-eighth the allocation cost.
		\item Learned importance coordinate gives a $1.19\times$ FOM ratio, bounded by a two-level bootstrap.
		\item At a depth where naive is provably blind, self-consistency training gives a significant reliability win: 66.7\% vs.\ 40.0\% ($p=0.017$), the paper's one significant positive result.
		\item Uniform exact/validated/open grading maps where the transport toolkit transfers to spin systems.
	\end{highlights}
	
	% ---- Keywords ----------------------------------------------------------
	\begin{keywords}
		Rare-event sampling \sep Adjoint transport \sep Committor function \sep Weighted Ensemble \sep FW-CADIS \sep Variance reduction \sep Spin systems
	\end{keywords}
	
	\maketitle
	
	% ========================================================================
	\section{Introduction}\label{sec:intro}
	
	Rare-event processes, including phase transitions, conformational changes, and chemical reactions, are central to many problems in physics, chemistry, and materials science. Their efficient simulation requires importance to be concentrated in regions that contribute most strongly to the event of interest. For an event with probability $p$, the relative statistical error of an unbiased Monte Carlo estimator scales as $(Np)^{-1/2}$, motivating importance-sampling strategies such as importance sampling and trajectory splitting \cite{kahn1951,kahn1953,allen2009,touchette2009}.
	
	In statistical mechanics, several rare-event methods exploit this principle. Transition Path Theory characterizes reactive trajectories through the committor function $h(s)$, defined as the probability of reaching a target state before returning to the initial state, and identifies it as an optimal reaction coordinate \cite{e2006,metzner2009,evanden2010}. Practical approaches include Weighted Ensemble, Forward Flux Sampling, Adaptive Multilevel Splitting, and Transition Path Sampling \cite{huber1996,zhang2010,allen2009,cerou2007,bolhuis2002}. These methods have been applied to a wide range of rare molecular and phase-transition processes, with their performance depending strongly on the choice of a suitable progress coordinate \cite{allen2009,zuckerman2017}.
	
	A closely related importance function is well established in radiation transport. The adjoint flux describes the expected contribution of a particle introduced at a given phase-space point to a detector response and can be used to construct importance maps for particle splitting and Russian roulette \cite{lewins1965,bell1970,booth1984}. The CADIS methodology provides a systematic framework for generating source biasing and weight windows from an adjoint solution, while FW-CADIS extends this approach to multiple responses by targeting a more uniform statistical accuracy \cite{wagner1998,haghighat2003,wagnerfw2014,cooper2001}. Performance is commonly quantified using the figure of merit,
	$\mathrm{FOM}=1/(\sigma_{\mathrm{rel}}^{2}T)$, which combines statistical precision and computational cost \cite{spanier1969,wagnerfw2014}.
	
	The connection between these two fields arises from the underlying mathematical structure of their importance functions. For suitable Markov processes, the adjoint function associated with an absorbing target and the rare-event committor satisfy corresponding harmonic equations. This relationship can be expressed through Doob's $h$-transform, which provides the optimal change of measure for conditioning a process on a rare event \cite{doob1984}. Related formulations appear in large-deviation theory and stochastic control, where the corresponding harmonic function defines an optimal or zero-variance importance transformation \cite{chetrite2015,touchette2009,hartmann2019}. This correspondence suggests that adjoint-based variance-reduction concepts can provide a systematic framework for constructing and evaluating committor-based importance functions in statistical-mechanical systems.
	
	Recent developments in machine learning provide practical approaches for estimating committors in high-dimensional systems. Physics-informed neural networks and related variational approaches have been used to solve committor equations directly \cite{raissi2019,li2019,khoo2019}, while active-learning strategies can focus training on dynamically important regions \cite{rotskoff2020}. Other approaches have employed learned committors as collective variables for rare-event sampling \cite{trizio2025}. Generative neural models have also demonstrated the ability to represent complex equilibrium distributions in spin systems \cite{wu2019,noe2019,nicoli2020,gabrie2022}. Together, these developments provide a basis for combining learned importance functions with established variance-reduction principles.
	
	This work has investigated such a combination by transferring the adjoint-based importance framework of CADIS/FW-CADIS to lattice spin systems using a learned committor as the importance function. The study focuses on three objectives: (i) establish the correspondence between the committor and the optimal importance function for a specified rare-event tally; (ii) investigate the use of CADIS/FW-CADIS-inspired weight-window strategies with an approximate learned importance function; and (iii) evaluate performance using the figure of merit as a cost-aware measure of variance reduction. The approach is examined using both the two-dimensional Ising ferromagnet \cite{ising1925,onsager1944} and the Edwards--Anderson model \cite{edwards1975,binder1986}, providing systems with different levels of structural complexity. Metropolis dynamics are used as the underlying transition mechanism \cite{metropolis1953}, while the resulting variance-reduction performance is assessed directly through rare-event tallies and computational efficiency.
	
	\FloatBarrier
	
	% ========================================================================
	
	
	\section{Theoretical Approach}\label{sec:methods}
	
	\subsection{Adjoint importance and CADIS}\label{sec:theory}
	
	Let $P=(\mathbf r,\mathbf\Omega,E)$ denote a phase-space point and let $\psi(P)$ be the particle flux satisfying the linear transport equation
	\begin{equation}
		\mathcal{H}\psi=q,
	\end{equation}
	where $\mathcal{H}$ is the transport operator and $q$ is the source. A detector response is a linear functional
	\begin{equation}
		R=\langle \sigma_d,\psi\rangle ,
	\end{equation}
	where $\sigma_d$ represents the detector response function. The corresponding adjoint flux $\psi^\dagger$ satisfies
	\begin{equation}
		\mathcal{H}^{\dagger}\psi^\dagger=\sigma_d .
		\label{eq:adjoint}
	\end{equation}
	Using the adjoint relationship,
	\begin{equation}
		R=\langle \sigma_d,\psi\rangle
		=\langle \psi^\dagger,q\rangle .
		\label{eq:fund}
	\end{equation}
	Thus, $\psi^\dagger(P)$ represents the expected contribution of a particle introduced at $P$ to the detector response and can therefore be interpreted as an importance function.
	
	The exact adjoint provides the theoretical zero-variance importance function. In practice, an approximate adjoint is used to bias particle transport and control statistical weights. CADIS (Consistent Adjoint-Driven Importance Sampling) constructs a biased source and weight-window parameters from an approximate adjoint solution \cite{wagner1998,haghighat2003}. The target particle weight is chosen approximately as
	\begin{equation}
		w(P)\propto \frac{1}{\psi^\dagger(P)},
	\end{equation}
	with particles split or subjected to Russian roulette when their weights move outside a prescribed window. FW-CADIS extends this framework to multiple responses by constructing a global importance distribution that promotes more uniform relative uncertainty across a family of tallies \cite{wagnerfw2014,cooper2001}. Efficiency is quantified using the figure of merit
	\begin{equation}
		\mathrm{FOM}=\frac{1}{\sigma_{\mathrm{rel}}^2T},
		\label{eq:fom}
	\end{equation}
	where $\sigma_{\mathrm{rel}}$ is the relative statistical error and $T$ is the computational time.
	
	\subsection{Adjoint--committor correspondence}\label{sec:identity}
	
	Consider a Markov process on configurations $s$ with transition kernel $K(s\rightarrow s')$. Let $A$ denote the initial or bulk set and $B$ the target set. The committor is defined as
	\begin{equation}
		h(s)=\Pr[\tau_B<\tau_A\mid s],
	\end{equation}
	with boundary conditions
	\begin{equation}
		h|_A=0,\qquad h|_B=1.
	\end{equation}
	On the transient states, the committor satisfies the backward equation
	\begin{equation}
		h(s)=\sum_{s'}K(s\rightarrow s')h(s'),
		\qquad\text{or}\qquad
		(\mathbb I-K)h=0 .
		\label{eq:committor}
	\end{equation}
	
	This equation has the same mathematical form as the discretized adjoint transport equation when the target set is treated as the response region. Consequently, the committor and adjoint importance function represent the same underlying quantity: the expected contribution of the current state to a specified future rare-event response. This correspondence is a direct consequence of the common Markov and linear-operator structure and is closely related to Doob's $h$-transform \cite{doob1984}. For the exact committor, the transformed transition kernel
	\begin{equation}
		\widetilde K(s\rightarrow s')
		=K(s\rightarrow s')\frac{h(s')}{h(s)}
	\end{equation}
	provides the zero-variance change of measure for the corresponding hitting probability \cite{chetrite2015,touchette2009}. The committor therefore provides a natural importance function for rare-event sampling, while approximate committors require a variance-control mechanism such as population splitting and resampling.
	
	For the present study, the target probability is the dynamical hitting probability from a reference bulk state. This quantity is distinct from the static equilibrium probability of finding the system inside the target set. Under the mixing conditions assumed in transition-path theory, the dynamical and equilibrium descriptions can be related through equilibration within the bulk set before escape \cite{e2006}.
	
	\subsection{Scope of the correspondence}\label{sec:scope}
	
	The correspondence is mathematical and applies to any Markov process with an absorbing target set. After discretization, a transport problem can be represented by a finite transition operator,
	\begin{equation}
		(\mathbb I-\mathbf K)\boldsymbol{\psi}^{\dagger}
		=\boldsymbol{\sigma}_d ,
	\end{equation}
	which has the same linear structure as the committor equation in Eq.~\eqref{eq:committor}. The physical interpretation of $\mathbf K$, however, differs between applications. In radiation transport it represents particle streaming, scattering, and absorption, whereas in spin systems it represents transitions between configurations generated by the chosen spin dynamics.
	
	This distinction does not affect the underlying variance-reduction principle. In the present work, the adjoint-committor correspondence is therefore used to transfer the concept of an importance function from radiation transport to lattice spin dynamics, while replacing the analytically or numerically computed adjoint with a learned approximation of the committor.
	
	\subsection{Spin models and rare-event targets}\label{sec:models}
	
	Two $L\times L$ spin systems with periodic boundaries are considered under the Gibbs distribution
	\begin{equation}
		p(s)\propto \exp[-E(s)/T].
	\end{equation}
	The first is the two-dimensional ferromagnetic Ising model,
	\begin{equation}
		E(s)=-J\sum_{\langle ij\rangle}s_i s_j,\qquad J>0,
	\end{equation}
	for which magnetization provides a natural reaction coordinate. The second is the Edwards--Anderson (EA) spin glass, in which the couplings $J_{ij}=\pm J$ are quenched and disordered \cite{edwards1975,binder1986}. The EA model provides a more complex configuration space in which magnetization alone does not generally provide a sufficient description of the relevant rare-event pathways.
	
	Unless otherwise stated, simulations use $T=2.6$, above the critical temperature of the two-dimensional Ising model, with $L=16$ for production calculations and $L=4$ for exact enumeration. Single-spin-flip Metropolis dynamics are used to generate the underlying Markov process \cite{metropolis1953}.
	
	Rare-event targets are defined as nested energy thresholds,
	\begin{equation}
		B_k=
		\left\{
		s:\frac{E(s)}{N}\leq\epsilon_k
		\right\},
	\end{equation}
	where
	\begin{equation}
		\epsilon_k=E_{\min}^{\rm pilot}-k\Delta,
		\qquad
		\Delta=\frac{\mu-E_{\min}^{\rm pilot}}{10}.
		\label{eq:beyond}
	\end{equation}
	Here, $\mu$ and $E_{\min}^{\rm pilot}$ are the mean and minimum energy per spin obtained from a preliminary simulation. The parameter $k$ therefore provides a systematic rarity scale, while the nested structure
	$B_1\supset B_2\supset\cdots$
	allows the same target family to be used for multi-threshold variance reduction.
	
	\subsection{Weighted Ensemble and learned importance}\label{sec:we}
	
	Weighted Ensemble (WE) provides the population-control mechanism used in this study. A population of weighted replicas is periodically assigned to bins according to a progress or importance coordinate. Replicas are split or merged to approach a prescribed occupancy while conserving the total statistical weight, preserving unbiasedness for arbitrary binning coordinates \cite{huber1996,zhang2010}. In the present framework, the learned importance function $I_\theta(s)$ replaces a manually selected reaction coordinate and is used to define the WE bins.
	
	Three approaches are considered for estimating $I_\theta(s)$. The first uses a surrogate rollout objective in which short trajectories from sampled configurations provide a continuous measure of the best progress toward the target. The second uses a milestoning formulation in which intermediate thresholds provide local estimates of progress toward successive milestones. The third follows the self-consistency approach of Kim and Cai (2026), training the network to satisfy
	\begin{equation}
		I_\theta(s)\approx\sum_jK(s\rightarrow j)I_\theta(j),
	\end{equation}
	with appropriate boundary conditions at the target and bulk states.
	
	These approaches provide increasingly direct approximations to the committor while avoiding the difficulty of obtaining sufficient direct labels for extremely rare target events.
	
	\subsection{FW-CADIS global allocation}\label{sec:fwcadis-method}
	
	For the nested target family, FW-CADIS-inspired allocation is used to distribute computational effort across the importance bins. Given an estimated equilibrium occupancy $P(b)$ for bin $b$, the number of replicas assigned to that bin is parameterized as
	\begin{equation}
		N_b\propto [P(b)]^{-\alpha},
	\end{equation}
	where $\alpha=0$ corresponds to uniform allocation and $\alpha=1$ corresponds to the full inverse-occupancy weighting. The two cases use the same sampling and propagation procedures and differ only in the allocation rule.
	
	Performance is evaluated first by target coverage, i.e., the number of thresholds reached by the simulations, and subsequently by the uniformity of relative statistical error across the observed thresholds. This evaluation is complemented by the figure of merit in Eq.~\eqref{eq:fom}, allowing both statistical precision and computational cost to be considered.
	
	
	\FloatBarrier
	
	
	% ========================================================================
	\section{Results and Discussion}\label{sec:results}
	
	\subsection{VAN and the committor demo}\label{sec:res-van}
	The VAN was trained at $L=16$ and reproduces the exact Onsager free energy across the
	tested temperature range. Its purpose in this project is to remove critical slowing down,
	not to sample rare events: a separate autocorrelation measurement at $L=32$ found the
	integrated autocorrelation time of ordinary Metropolis sampling spikes to $\tau_{\rm
		int}=46.1$ sweeps at $T_c\approx2.269$, versus 1.5--3 sweeps away from criticality, while
	the VAN samples independently by construction and pays none of this cost. This is a speed
	result, orthogonal to the rare-event problem the rest of the paper addresses (they are
	different difficulties, \S\ref{sec:methods}).
	
	The committor identity of \S\ref{sec:identity} was checked on a solvable 1D absorbing
	random walk on $\{0,\dots,15\}$ before any lattice code was written, with the rare event
	(reaching state 15 before 0) calibrated to $\pi=1.14\times10^{-3}$. Three samplers were
	compared. Against naive Monte Carlo as the baseline, tilting with the exact committor gave
	a measured relative variance of ${\sim}10^{-30}$, zero to machine precision, as the
	zero-variance argument of \S\ref{sec:identity} requires. Tilting with an
	\emph{approximate} committor under pure biasing, with no weight windows, was
	\emph{worse than naive}.
	
	That third outcome is the reason every subsequent learned-coordinate result in this paper
	uses Weighted Ensemble (split/roulette) rather than raw importance reweighting: an
	imperfect importance estimate, used to reweight without bounding the weights, can be
	actively harmful.
	
	\subsection{Weighted Ensemble: validated}\label{sec:res-we}
	At $L=4$, WE reproduces the exact enumerated $P(m)$ to a mean $|\Delta\ln P|$ of
	0.006--0.024 across three tested regimes. At $L=20$, WE reaches roughly two decades deeper
	into the magnetization tail than naive Monte Carlo at matched cost: naive's hard floor is
	$|m|=0.886$ ($P\approx6\times10^{-6}$), where it returns zero samples and hence zero
	information; WE reaches $|m|=0.932$ ($P\approx2.4\times10^{-7}$). Separately, applying
	weight windows to the same \emph{approximate}, otherwise-fragile tilt of
	\S\ref{sec:res-van} recovered roughly a $2\times$ figure-of-merit improvement
	over the same tilt without them (measured FOM $1.24\times10^{-5}\to4.86\times10^{-4}$):
	the robustness layer converts an actively harmful bias into a net improvement, exactly as
	the transport literature predicts. Magnetization is confirmed useless as a WE coordinate
	for the EA spin glass specifically: an early test recorded 10/10 zero-replica failure for
	$\mathrm{WE}[m]$ where energy-binning still found the target, the direct empirical
	confirmation of the frustration argument in \S\ref{sec:methods}.
	
	\subsection{Learned coordinate vs.\ hand-picked coordinates}\label{sec:res-learned}
	Table~\ref{tab:prefix} reports the full-scale ($L=16$, 10 replicas) result on both models,
	at a target calibrated one step beyond the pilot's observed extreme.
	
	
	Two readings, both necessary (Table~\ref{tab:prefix}): naive Monte Carlo wins on raw FOM in both models at this
	rarity ($\sim\!4\times10^{-5}$), and on the spin glass, with $\mathrm{WE}[E]$ and
	$\mathrm{WE}[I_\theta]$ on equal footing (both 2/10 zero-replicas), the learned
	coordinate has the higher FOM by $\sim$27\% ($6.36\times10^{-9}$ vs.\ $5.00\times10^{-9}$),
	the first apples-to-apples win for the learned coordinate over the best hand-picked
	one, and the direct predecessor of the milestoning result below. $\mathrm{WE}[m]$ fails in
	most spin-glass replicas (7/10), the direct empirical confirmation of the frustration
	argument in \S\ref{sec:methods}: magnetization carries essentially no information about
	this landscape.
	% --- TABLE 1 ---
	\begin{table}[H]
		\centering\small
		\caption{Full scale, both models. FOM vs.\ $\mathrm{WE}[m]$ in parentheses.}
		\label{tab:prefix}
		\setlength{\tabcolsep}{4pt}
		\begin{tabular}{lrrrrr}
			\toprule
			Model / Method & $\hat\pi$ & rel.\ s.d.\ & cost & FOM (vs.\ $\mathrm{WE}[m]$) & zero-rep.\ \\
			\midrule
			\multicolumn{6}{l}{\emph{Ferromagnet}, target $|m|\ge0.9297$} \\
			naive & $9.52\times10^{-6}$ & 0.441 & $1.34\times10^{8}$ & $3.83\times10^{-8}$ ($2.92\times$) & 0/10 \\
			$\mathrm{WE}[m]$ & $1.05\times10^{-5}$ & 0.614 & $2.02\times10^{8}$ & $1.31\times10^{-8}$ ($1.00\times$) & 0/10 \\
			$\mathrm{WE}[E]$ & $1.09\times10^{-5}$ & 1.343 & $9.81\times10^{7}$ & $5.65\times10^{-9}$ ($0.43\times$) & 2/10 \\
			$\mathrm{WE}[I_\theta]$ & $7.92\times10^{-6}$ & 0.893 & $1.82\times10^{8}$ & $6.90\times10^{-9}$ ($0.53\times$) & 0/10 \\
			\multicolumn{6}{l}{\emph{EA spin glass}, target $E/N\le-1.0625$} \\
			naive & $1.00\times10^{-5}$ & 0.633 & $1.34\times10^{8}$ & $1.86\times10^{-8}$ ($2.38\times$) & 1/10 \\
			$\mathrm{WE}[m]$ & $1.02\times10^{-5}$ & 1.857 & $3.72\times10^{7}$ & $7.80\times10^{-9}$ ($1.00\times$) & 7/10 \\
			$\mathrm{WE}[E]$ & $1.65\times10^{-5}$ & 1.115 & $1.61\times10^{8}$ & $5.00\times10^{-9}$ ($0.64\times$) & 2/10 \\
			$\mathrm{WE}[I_\theta]$ & $1.54\times10^{-5}$ & 1.052 & $1.42\times10^{8}$ & $6.36\times10^{-9}$ ($0.81\times$) & 2/10 \\
			\bottomrule
		\end{tabular}
	\end{table}
	
	Is $I_\theta$ actually informative, or just a noisy copy of energy? A separate
	diagnostic regresses the rare-event label on energy alone, on $I_\theta$ alone, and on
	both, at matched training budget (Table~\ref{tab:diagnostic}).
	
	% --- TABLE 2: diagnostic regression (relocated from the appendix) ---
	\begin{table}[H]
		\centering\small
		\caption{Incremental-$R^2$ diagnostic, EA spin glass, $L=16$, matched training budget
			(MSE $1.0\to0.271$ over 20 epochs).}
		\label{tab:diagnostic}
		\begin{tabular}{lr}
			\toprule
			Quantity & value \\
			\midrule
			Pearson corr$(I_\theta,\,E/N)$ & $-0.614$ \\
			fraction of $\mathrm{Var}(I_\theta)$ unexplained by $E/N$ & $0.623$ \\
			$R^2(\text{label}\sim E/N\text{ alone})$ & $0.358$ \\
			$R^2(\text{label}\sim I_\theta\text{ alone})$ & $0.577$ \\
			$R^2(\text{label}\sim E/N+I_\theta)$ & $0.608$ \\
			incremental $R^2$ from $I_\theta$ beyond $E/N$ & $+0.250$ \\
			\bottomrule
		\end{tabular}
	\end{table}
	
	$+0.250$ is far above the diagnostic's own noise threshold ($0.02$): $I_\theta$ carries
	substantial genuine information about the spin glass that energy alone does not have,
	consistent with the FOM advantage measured directly in Table~\ref{tab:prefix}.
	
	Pushing rarer does not (yet) resolve the naive-vs-learned question. Rerunning one
	threshold step rarer (``beyond=2'') leaves the comparison uninterpretable rather than
	resolving it in either direction: 4--10 of 10 replicas fail across every method at this
	depth (naive 4/10, $\mathrm{WE}[E]$ 5/10, $\mathrm{WE}[I_\theta]$ 7/10, and
	$\mathrm{WE}[m]$ 10/10, i.e.\ no estimate at all), so the numbers are noise over noise,
	not a result. This
	is the run that established the ladder-difficulty discipline of \S\ref{sec:methods} for the
	remainder of the project.
	
	Notably, \emph{naive itself} now has 4/10 zero-replicas (versus 1/10 at the shallower
	target). This rarity is stretching every method's budget, not only the weight-window
	ones, so a trustworthy comparison here needs a different fix than anything tried at the
	shallower depth. This motivated the milestoning objective below, which is a change to the
	training label rather than to the weight-window schedule.
	
	\subsection{Milestoning: a clean regime for comparison, but wildly seed-dependent}\label{sec:milestone}
	On the EA spin glass at $L=16$, one threshold step beyond the pilot's observed extreme
	(``beyond=1''), every method reaches 0/10 zero-replicas, the cleanest,
	most-trustworthy regime obtained in this project by the usual ladder-difficulty standard
	(Table~\ref{tab:milestone}).
	
	% --- TABLE 3: 17-seed milestoning sweep (relocated from the appendix) ---
	\begin{table}[H]
		\centering\footnotesize
		\caption{EA spin glass, $L=16$, $T=2.6$, one step beyond the pilot extreme, 10 replicas,
			milestoning labels, seventeen independent training runs of the identical configuration.
			Every method observed the event in every replica in every run; naive,
			$\mathrm{WE}[m]$, and $\mathrm{WE}[E]$ are bit-identical across all runs (deterministic,
			and $\mathrm{WE}[E]$ does not depend on the trained network). Seed 0 was run twice
			independently and reproduced bit-for-bit, confirming that the seeding fix described in
			the text is complete.}
		\label{tab:milestone}
		\setlength{\tabcolsep}{4pt}
		\begin{tabular}{lrrrr}
			\toprule
			Method & $\hat\pi$ & rel.\ s.d.\ & FOM & zero-replicas \\
			\midrule
			naive & $5.29\times10^{-5}$ & 0.225 & $1.48\times10^{-7}$ & 0/10 \\
			$\mathrm{WE}[m]$ & $5.55\times10^{-5}$ & 0.201 & $1.17\times10^{-7}$ & 0/10 \\
			$\mathrm{WE}[E]$ & $4.49\times10^{-5}$ & 0.209 & $4.22\times10^{-8}$ & 0/10 \\
			$\mathrm{WE}[I_\theta]$, unseeded run 1 & $5.03\times10^{-5}$ & 0.174 &
			$4.81\times10^{-8}$ ($+14\%$) & 0/10 \\
			$\mathrm{WE}[I_\theta]$, unseeded run 2 & $4.52\times10^{-5}$ & 0.204 &
			$3.48\times10^{-8}$ ($-17\%$) & 0/10 \\
			$\mathrm{WE}[I_\theta]$, seed 0 (runs 3a/3b) & $5.28\times10^{-5}$ & 0.127 &
			$8.94\times10^{-8}$ ($+112\%$) & 0/10 \\
			$\mathrm{WE}[I_\theta]$, seed 1 & $5.09\times10^{-5}$ & 0.216 &
			$3.09\times10^{-8}$ ($-27\%$) & 0/10 \\
			$\mathrm{WE}[I_\theta]$, seed 2 & $5.06\times10^{-5}$ & 0.175 &
			$4.65\times10^{-8}$ ($+10\%$) & 0/10 \\
			$\mathrm{WE}[I_\theta]$, seed 3 & $4.56\times10^{-5}$ & 0.252 &
			$2.29\times10^{-8}$ ($-46\%$) & 0/10 \\
			$\mathrm{WE}[I_\theta]$, seed 4 & $4.13\times10^{-5}$ & 0.128 &
			$8.81\times10^{-8}$ ($+109\%$) & 0/10 \\
			$\mathrm{WE}[I_\theta]$, seed 5 & $5.40\times10^{-5}$ & 0.138 &
			$7.65\times10^{-8}$ ($+82\%$) & 0/10 \\
			$\mathrm{WE}[I_\theta]$, seed 6 & $4.69\times10^{-5}$ & 0.149 &
			$6.56\times10^{-8}$ ($+56\%$) & 0/10 \\
			$\mathrm{WE}[I_\theta]$, seed 7 & $4.32\times10^{-5}$ & 0.267 &
			$2.03\times10^{-8}$ ($-52\%$) & 0/10 \\
			$\mathrm{WE}[I_\theta]$, seed 8 & $4.89\times10^{-5}$ & 0.135 &
			$7.86\times10^{-8}$ ($+86\%$) & 0/10 \\
			$\mathrm{WE}[I_\theta]$, seed 9 & $5.03\times10^{-5}$ & 0.159 &
			$5.66\times10^{-8}$ ($+34\%$) & 0/10 \\
			$\mathrm{WE}[I_\theta]$, seed 10 & $4.80\times10^{-5}$ & 0.157 &
			$5.85\times10^{-8}$ ($+39\%$) & 0/10 \\
			$\mathrm{WE}[I_\theta]$, seed 11 & $4.81\times10^{-5}$ & 0.202 &
			$3.48\times10^{-8}$ ($-17\%$) & 0/10 \\
			$\mathrm{WE}[I_\theta]$, seed 12 & $5.14\times10^{-5}$ & 0.230 &
			$2.76\times10^{-8}$ ($-35\%$) & 0/10 \\
			$\mathrm{WE}[I_\theta]$, seed 13 & $4.83\times10^{-5}$ & 0.150 &
			$6.36\times10^{-8}$ ($+51\%$) & 0/10 \\
			$\mathrm{WE}[I_\theta]$, seed 14 & $4.94\times10^{-5}$ & 0.146 &
			$6.76\times10^{-8}$ ($+60\%$) & 0/10 \\
			$\mathrm{WE}[I_\theta]$, seed 15 & $4.53\times10^{-5}$ & 0.232 &
			$2.66\times10^{-8}$ ($-37\%$) & 0/10 \\
			\bottomrule
		\end{tabular}
	\end{table}
	
	Every $\mathrm{WE}[I_\theta]$ row of Table~\ref{tab:milestone} is, as far as the pipeline
	is concerned, the
	\emph{same experiment}: identical target, identical training data (from cached labels),
	identical hyperparameters. The only thing that varies is the network's random weight
	initialization (\texttt{torch.manual\_seed} was absent from this driver; found and
	fixed during this investigation, then swept explicitly via \texttt{--net-seed}).
	Seventeen seeds give seventeen different outcomes, ranging from a $52\%$ loss to
	a $112\%$ win against $\mathrm{WE}[E]$, a wide, roughly symmetric spread around a
	small positive average. The seeding fix itself is confirmed complete and correct: seed
	0 was run twice independently (runs 3a and 3b) and reproduced bit-for-bit, including the
	full 1774--1776s wall-clock trajectory, ruling out any other unseeded source of
	variation in this pipeline. Per-replica values needed for bootstrapping were retained
	for seventeen of the eighteen runs attempted (all but unseeded run 1, which predates that
	logging).
	Pooling those seventeen with a two-level bootstrap, outer resample over seeds,
	inner resample over each seed's 10 replicas, the same discipline used throughout
	\S\ref{sec:fwcadis} and \S\ref{sec:kimcai} but extended with the seed axis, gives a
	geometric-mean FOM ratio of $1.19\times$ with 95\% CI $[0.84, 1.75]$ (median-of-ratios
	agrees: $1.20\times$, CI $[0.74, 1.85]$; the arithmetic mean of ratios is not used here
	because it is unstable at $R=10$ replicas per seed: some bootstrap draws land on
	near-duplicate resamples with rel.\ s.d.\ near zero, and since $\mathrm{FOM}\propto
	1/\mathrm{rel.\,var.}$ that blows up into large outliers that skew a plain mean).
	The estimate was tracked at four checkpoints as seeds were added: $1.28\times$
	(n=7), $1.17\times$ (n=9), $1.28\times$ (n=12), $1.19\times$ (n=17), and it does
	\emph{not} trend in either direction. What does move monotonically is the CI width
	($1.73\to1.36\to1.18\to0.91$), the expected signature of averaging down noise around a
	small effect whose sign is not yet resolved, not evidence of a shrinking or growing true
	effect. We report the full trajectory rather than the final number alone, because citing
	only the endpoint would silently launder a mid-sequence checkpoint (n=9's ``moving
	toward the null'') that did not survive more data, a mistake this project's own
	methodology is built to catch.
	Extrapolating the observed CI-width-vs-$n$ scaling against the current point
	estimate, reaching a lower bound that excludes 1 would need roughly 100 total seeds
	($\sim\!40$ more hours of compute), so this was not pursued further. This is reported as
	the honest terminus of the comparison, not an open
	thread: seventeen seeds is already an unusually thorough characterisation, the point
	estimate is stably positive but modest ($\sim\!1.2\times$) across four independent
	checkpoints (Figure~\ref{fig:seedsweep}), and the CI cannot be closed by seed count alone
	at a cost proportionate to
	the question. Naive Monte Carlo has the best raw FOM in every run
	at this specific rarity ($\sim\!5\times10^{-5}$) regardless of seed. Variance
	reduction has not yet earned back its overhead here, the same lesson demonstrated in
	miniature by the 1D toy-walk result in \S\ref{sec:res-van}.
	
	% --- FIGURE 4 ---
	\noindent
	\begin{minipage}{\linewidth}
		\centering
		\includegraphics[width=0.92\textwidth]{figs/fig_seed_sweep.pdf}
		\captionof{figure}{Seed-to-seed variability of the learned milestoning coordinate.
			Left: all 17 per-seed FOM ratios of Table~\ref{tab:milestone} (sorted), with the pooled
			two-level-bootstrap
			median and 95\% CI shaded. Right: the four tracked checkpoints as seeds were added,
			the point estimate bounces while the CI width shrinks monotonically, the expected
			signature of averaging down noise around a small effect whose sign is not resolved by
			seed count alone.}
		\label{fig:seedsweep}
	\end{minipage}
	\par\vspace{0.8\baselineskip}
	
	A plausible infrastructure explanation was tested and ruled out. One candidate
	cause for the seed-to-seed spread is that $\mathrm{WE}[I_\theta]$'s bin range was left
	adaptive (auto-ranging from the population actually explored) rather than calibrated to
	the coordinate's own scale, unlike $\mathrm{WE}[E]$ and $\mathrm{WE}[m]$. The diagnostic
	regression of Table~\ref{tab:diagnostic}
	(\texttt{05\_neural\_importance/diagnose\_coordinate.py}) had
	already shown $I_\theta$ carries real predictive signal beyond energy (incremental
	$R^2=+0.25$), with the verdict: if it still does not beat $\mathrm{WE}[E]$, look at the
	schedule, not the map. Two calibrations were tried: a theoretical one from the label's
	own z-score normalization, verified \emph{before} spending compute on it to be wrong (the
	network's real output never approaches those extremes, since training MSE never reaches
	0); and an empirical one, from what each trained network actually outputs on the pilot
	population and training set. The empirical calibration was run head-to-head against the
	untouched adaptive baseline on matched net-seeds (identical trained network per pair,
	only the bin range differing): mean FOM ratio \emph{worsened} ($0.79\times\to0.61\times$)
	and rel.\ s.d.\ \emph{worsened} ($0.214\to0.286$) across seeds 1--3. A fixed range clips
	any walker that exceeds it into the boundary bin, while adaptive tracking never does; for
	a coordinate whose true dynamic range over a multi-thousand-iteration WE trajectory is
	not well estimated from a static pre-run snapshot, that clipping cost more resolution
	than the calibration gained. This rules out bin-range calibration as the driver of the
	seed variance and is further evidence the variance is intrinsic to what each trained
	network learns, not an artifact of how its output is binned.
	
	Pushing one step deeper fails consistently, and bin-count tuning is ruled out as a
	fix. Every configuration tried at one threshold step rarer (``beyond=2'') was scored:
	the default bin count ($n_{\rm bins}=50$) at two replica counts, and a $1.8\times$ larger
	bin count. $\mathrm{WE}[I_\theta]$ scores \emph{worse} than $\mathrm{WE}[E]$ in all
	three. At the default configuration ($n_{\rm bins}=50$, $R=10$), where every method
	reaches 0/10 zero-replicas, the FOMs are naive $3.51\times10^{-8}$, $\mathrm{WE}[E]$
	$9.21\times10^{-9}$, $\mathrm{WE}[I_\theta]$ $8.36\times10^{-9}$, and $\mathrm{WE}[m]$
	$7.73\times10^{-9}$: naive ahead by roughly $4\times$, and the learned coordinate below
	the hand-picked one.
	
	Two further configurations were tried in an attempt to fix this with bin resolution instead
	of accepting it as structural: doubling the replica count ($R=10\to20$, same $n_{\rm
		bins}=50$) reproduced the same qualitative pattern ($\mathrm{WE}[I_\theta]$ still below
	$\mathrm{WE}[E]$); and raising the bin count by $1.8\times$ ($n_{\rm bins}=50\to90$, $R=10$)
	made $\mathrm{WE}[I_\theta]$'s \emph{own} FOM \emph{worse}, not better
	($8.36\times10^{-9}\to5.36\times10^{-9}$), which rules out bin resolution as the fix, since
	a genuine resolution shortfall would improve with more bins rather than degrade. The
	mechanism is structural: the learned coordinate's walker-population cost scales with bin
	count, while a discretized coordinate like energy naturally saturates a small, fixed
	number of bins and does not pay that overhead as it grows.
	
	A second lever, the WE resampling budget itself, was also tried
	($\mathrm{we\_iter}$: $1000\to2000$, same $n_{\rm bins}=50$) and produced a similarly
	unfavorable result: every weight-window method's FOM got \emph{worse}, not better, and
	$\mathrm{WE}[I_\theta]$ worst of all. That specific run used $\tau=3$ rather than this
	section's $\tau=2$, confounding the budget increase with a change in resampling interval,
	and a matched rerun of it specifically was not completed. The broader question, whether
	more WE budget helps, was instead answered more thoroughly by the deep-rare
	budget-scaling test of \S\ref{sec:deeprare}: scaling the budget up does not close the gap
	against naive at a shallow threshold depth either, consistent with, not contradicting,
	the bin-resolution finding above.
	
	The resampling interval between split/merge operations matters: it must be comparable to
	the system's own correlation time ($\sim$1 sweep at $T=2.6$), or split walkers fully relax
	and forget their split origin before the next resampling, degenerating every
	$\mathrm{WE}[\cdot]$ variant to naive Monte Carlo with bookkeeping overhead. Every
	milestoning result in this section uses an interval of 2 sweeps.
	
	\subsection{The Kim \& Cai head-to-head}\label{sec:kimcai}
	The self-consistency training objective of \S\ref{sec:we} was implemented, gated
	against exact enumeration at $L=4$ (learned-coordinate unbiasedness ratio 1.004), and run
	head-to-head against the surrogate objective, first at laptop scale, then at full scale.
	
	First evidence, at laptop scale ($L=12$). At the shallowest calibrated target, the
	self-consistency objective's point estimate already exceeded the surrogate's
	($4.98\times10^{-8}$ vs.\ $2.93\times10^{-8}$, $\sim$70\% higher FOM, fewer zero-replicas:
	2/6 vs.\ 3/6), though it still fell short of $\mathrm{WE}[E]$'s $7.53\times10^{-8}$ at this
	rarity. Training cost is real: $\sim$94$\times$ longer than the surrogate objective at this
	scale ($\sim$752s vs.\ $\sim$8s), from the additional $N=L^2$ network evaluations the
	self-consistency loss requires per training configuration. This motivated running the
	comparison at full scale.
	
	The result, at full scale. At $L=16$, the shallowest calibrated target (``beyond''\,$=0$,
	the easiest regime in which every method reaches 0--1 zero-replicas out of 10), the
	self-consistency coordinate's \emph{point estimate} exceeds both the hand-picked energy
	coordinate and this project's own surrogate coordinate (Table~\ref{tab:kimcai},
	Appendix~\ref{app:tables}).
	
	This point estimate does not survive a bootstrap check, and the effect shrinks
	under replication. Both configurations were re-run with per-replica values retained, and
	the FOM ratio between methods bootstrapped by resampling the replica axis independently for
	each side (no shared random-number seed to pair on: uniform-style and biased allocations
	consume different numbers of walkers per bin, so seed-sharing between them desynchronises
	almost immediately, the same obstruction encountered independently in the FW-CADIS
	common-random-numbers attempt of \S\ref{sec:fwcadis}). Since training here is a fixed cost
	independent of replica count (a deterministic seed; only the cheap evaluation step scales
	with $R$), tripling to $R=30$ cost roughly six additional minutes on top of the
	$\sim$76-minute training already paid, cheap enough to just check
	(Table~\ref{tab:kimcaiboot}, first two rows of each block).
	
	% --- TABLE: merged Kim & Cai bootstrap (both depths) ---
	\begin{table}[H]
		\centering\small
		\caption{Bootstrapped FOM ratios, full scale ($L=16$), for the self-consistency objective
			against this project's surrogate and against the hand-picked energy coordinate. Ratio
			$>1$ favors self-consistency. The shallowest calibrated target
			($\mathrm{beyond}=0$) is discussed first; the $\mathrm{beyond}=1$ rows are taken up
			below. No interval excludes 1.}
		\label{tab:kimcaiboot}
		\begin{tabular}{lrrrr}
			\toprule
			Comparison & beyond & $R$ & FOM ratio (median) & 95\% CI \\
			\midrule
			self-consistent / surrogate & 0 & 10 & 1.83 & $[0.49,\,9.23]$ \\
			self-consistent / surrogate & 0 & 30 & 1.39 & $[0.65,\,2.95]$ \\
			self-consistent / surrogate & 1 & 10 & 1.14 & $[0.27,\,4.81]$\textsuperscript{a} \\
			self-consistent / $\mathrm{WE}[E]$ & 0 & 10 & 1.19 & $[0.33,\,4.95]$ \\
			self-consistent / $\mathrm{WE}[E]$ & 0 & 30 & 1.07 & $[0.55,\,2.09]$ \\
			self-consistent / $\mathrm{WE}[E]$ & 1 & 10 & 1.33 & $[0.34,\,5.74]$ \\
			\bottomrule
		\end{tabular}
		\\[2pt]
		{\footnotesize\textsuperscript{a}Zero-replicas at this depth: 2/10 for the surrogate
			objective against 1/10 for self-consistency.}
	\end{table}
	
	At $R=10$ neither ratio is significant. At $R=30$, both confidence intervals tighten
	substantially \emph{and} both median ratios move toward 1: not merely a tightening
	confidence interval around the original estimate, but the point estimate itself receding.
	We read the direction of this shift, not only its final inconclusiveness, as evidence that
	the original $R=10$ advantage was substantially noise, not signal. A network-caching
	mechanism was added after this result specifically to make a further replica increase (were
	one wanted) essentially free: training is deterministic, so only the cheap evaluation step
	needs to re-run to extend $R$ further, rather than repaying the fixed training cost again.
	
	Does the advantage grow at greater depth, where the surrogate objective's known
	weaknesses should bite harder? We tested the natural next question directly: rerunning
	both objectives one threshold step deeper (``beyond''\,$=1$, the same depth as
	Table~\ref{tab:milestone}); those are the $\mathrm{beyond}=1$ rows of
	Table~\ref{tab:kimcaiboot}.
	
	Not significant, as expected at $R=10$, but the \emph{direction} of the change from
	the $\mathrm{beyond}=0$ rows is itself informative and runs counter to the hypothesis being
	tested: the point-estimate gap versus the surrogate \emph{shrank} at this deeper target
	($1.83\times\to1.14\times$) rather than growing, even as the self-consistency coordinate's
	relative \emph{reliability} advantage held (fewest zero-replicas of any weight-window
	method at both depths tested). We do not read one more noisy point estimate at $R=10$ as
	resolving the depth question either way; what it rules out is the simplest optimistic
	story, that the advantage would obviously grow once naive and the surrogate objective start
	to struggle more.
	
	% --- FIGURE 5 ---
	\noindent
	\begin{minipage}{\linewidth}
		\centering
		\includegraphics[width=0.85\textwidth]{figs/fig_kimcai.pdf}
		\captionof{figure}{All six FOM-ratio comparisons of Table~\ref{tab:kimcaiboot} as a
			forest plot. Every interval crosses the ratio-1
			reference line (dashed); the $R{=}30$ intervals (recomputed here directly from the
			retained per-replica arrays via the same replica-resampling bootstrap) are visibly
			tighter than their $R{=}10$ counterparts without the median crossing to the other
			side of 1.}
		\label{fig:kimcai}
	\end{minipage}
	\par\vspace{0.8\baselineskip}
	
	Figure~\ref{fig:kimcai} collects all six comparisons tried.
	
	\subsection{Self-consistency in the deep-rare regime: the paper's one significant
		positive result}
	\label{sec:deeprare}
	Every comparison above (milestoning, Kim \& Cai at
	$\mathrm{beyond}=0/1$) was run at a depth where naive Monte Carlo remains fully
	reliable (0/10 zero-replicas). The theory this paper transfers has nothing to prove in
	that regime; its own domain (radiation shielding) is precisely where naive cannot reach
	at all. We calibrated where naive genuinely breaks down, at the same $L=16$, $T=2.6$
	configuration used throughout, using the custom-preset naive budget
	($\mathrm{naive\_chains}=150$, $\mathrm{naive\_sweeps}=3500$; Table~\ref{tab:deeprare}):
	
	\begin{table}[H]
		\centering\small
		\caption{Naive-only depth calibration, 10 replicas, custom-preset budget. Naive is
			fully reliable through $\mathrm{beyond}=2$, starts failing at $\mathrm{beyond}=3$,
			and is completely blind at $\mathrm{beyond}=4$ and beyond.}
		\label{tab:deeprare}
		\begin{tabular}{lrrr}
			\toprule
			beyond & $\hat\pi$ & rel.\ s.d.\ & zero-replicas \\
			\midrule
			1 & $5.29\times10^{-5}$ & 0.23 & 0/10 \\
			2 & $5.48\times10^{-6}$ & 0.46 & 0/10 \\
			3 & $7.14\times10^{-7}$ & 1.61 & 7/10 \\
			4 & $0.0$ & --- & 10/10 \\
			\bottomrule
		\end{tabular}
	\end{table}
	
	Simply targeting a deeper threshold was not, by itself, sufficient: at $\mathrm{beyond}=3$,
	scaling the WE budget up ($n_{\rm bins}:50\to80$, $\mathrm{we\_iter}:1000\to3000$) made
	$\mathrm{WE}[E]$'s FOM \emph{disadvantage against naive widen}, not narrow, since the
	extra reliability it bought cost more than it saved. Only at $\mathrm{beyond}=4$, where
	naive is not merely noisy but categorically blind, did the scaled budget produce a real
	capability gap: $\mathrm{WE}[E]$ found the event in 12 of 30 replicas where naive found
	essentially none (naive: 1/30, driven by a single lucky hit with rel.\ s.d.\ $=5.5$, not
	a meaningful measurement). $\mathrm{beyond}=4$ (threshold $E/N=-1.119$) with this scaled
	budget ($n_{\rm bins}=80$, $n_{\rm per\ bin}=40$, $\mathrm{we\_iter}=3000$, $\tau=2$) is
	the deepest, most naive-adverse regime tested anywhere in this paper, and the first
	place the self-consistency objective (\S\ref{sec:we}) was tested outside the
	shallow regime where naive wins anyway (Table~\ref{tab:selfconsistdeep}).
	
	\begin{table}[H]
		\centering\small
		\caption{EA spin glass, $L=16$, $\mathrm{beyond}=4$, self-consistency objective,
			three independent training seeds, 30 replicas each. $\mathrm{WE}[E]$/$\mathrm{WE}[m]$/naive
			do not depend on the trained network and are bit-identical across all three runs.}
		\label{tab:selfconsistdeep}
		\begin{tabular}{lrrr}
			\toprule
			Method & hits/30 & FOM & ratio vs.\ $\mathrm{WE}[E]$ \\
			\midrule
			naive & 1/30 & $2.48\times10^{-10}$\textsuperscript{a} & --- \\
			$\mathrm{WE}[m]$ & 5/30 & $1.04\times10^{-10}$ & --- \\
			$\mathrm{WE}[E]$ & 12/30 & $1.69\times10^{-10}$ & 1.00 \\
			$\mathrm{WE}[I_\theta]$, seed 0 & 19/30 & $2.18\times10^{-10}$ & $1.29\times$ \\
			$\mathrm{WE}[I_\theta]$, seed 1 & 20/30 & $1.81\times10^{-10}$ & $1.07\times$ \\
			$\mathrm{WE}[I_\theta]$, seed 2 & 21/30 & $1.77\times10^{-10}$ & $1.04\times$ \\
			\bottomrule
		\end{tabular}
		\\[2pt]
		{\footnotesize\textsuperscript{a}Not a meaningful FOM: driven by a single hit in 30
			replicas (rel.\ s.d.\ $=5.5$); reported for completeness only.}
	\end{table}
	
	Reliability: statistically significant. Pooled across the three independent
	networks, $\mathrm{WE}[I_\theta]$ found the event in $60/90$ replicas ($66.7\%$) versus
	$\mathrm{WE}[E]$'s $12/30$ ($40.0\%$); a Fisher exact test on this $2\times2$ table gives
	$p=0.017$. This is the first statistically significant \emph{positive} result for a
	learned coordinate anywhere in this paper. The three per-network rates ($63\%, 67\%,
	70\%$) are close enough to each other, unlike the milestoning coordinate's wide
	net-to-net FOM swings (\S\ref{sec:milestone}), that pooling them as a single binomial
	comparison is a reasonable approximation rather than one masking real heterogeneity.
	Figure~\ref{fig:deeprare} plots that comparison alongside the depth calibration of
	Table~\ref{tab:deeprare}, which is what makes it meaningful: the gap opens only outside
	the band in which every other comparison in this paper was run.
	
	% --- FIGURE: deep-rare reliability ---
	\noindent
	\begin{minipage}{\linewidth}
		\centering
		\includegraphics[width=0.92\textwidth]{figs/fig_deeprare.pdf}
		\captionof{figure}{The paper's one statistically significant positive result, and the
			depth at which it was obtained. Left: fraction of 30 replicas in which each method
			observed the rare event at $\mathrm{beyond}=4$
			(Table~\ref{tab:selfconsistdeep}), with Wilson 95\% intervals; white marks on the
			pooled bar are the three individual networks (19/30, 20/30, 21/30). Right: the
			naive-only depth calibration of Table~\ref{tab:deeprare}. Naive Monte Carlo is fully
			reliable through $\mathrm{beyond}=2$, the band in which every other comparison in
			this paper was run, and categorically blind at $\mathrm{beyond}=4$, where the
			self-consistency coordinate's advantage appears.}
		\label{fig:deeprare}
	\end{minipage}
	\par\vspace{0.8\baselineskip}
	
	FOM: directionally consistent, not yet significant. All three seeds' FOM
	ratios sit above 1 ($1.29\times, 1.07\times, 1.04\times$) with none of the sign-flipping
	seen everywhere else a coordinate was swept across seeds in this paper. A two-level
	bootstrap (outer over the 3 networks, inner over each network's 30 replicas) on the
	geometric-mean ratio gives $1.16\times$, 95\% CI $[0.65, 2.12]$, favorable but not
	excluding 1, since the underlying $\hat\pi$ estimates remain individually noisy
	(rel.\ s.d.\ $1.2$--$1.9$) even at 30 replicas in this regime.
	
	Reading this against \S\ref{sec:kimcai}'s claim that the self-consistency
	coordinate ``shows no sign of strengthening with depth'': that conclusion was drawn
	entirely within $\mathrm{beyond}\in\{0,1\}$, where naive remains competitive throughout.
	Tested at a depth where naive genuinely fails, the picture is materially different,
	consistent with the mechanistic reading that a committor-shaped training signal (the
	self-consistency objective enforces the harmonic condition $(I-K)h=0$ directly) carries
	information the milestoning surrogate cannot, since that surrogate's label is provably
	energy-adjacent (\texttt{progress\_value} for the spin glass is exactly $-E/N$; its
	rollout label is literally ``best energy reached in a short walk''), and two
	configurations at identical energy can have very different true committor values on a
	frustrated landscape. This is the paper's clearest evidence that the transport toolkit's
	distinctive machinery earns its keep specifically where naive cannot compete at all, not
	in the shallower regimes most of this paper's other comparisons were run in.
	
	\subsection{Does naive dominance persist at larger $L$?}\label{sec:lscaling}
	\S\ref{sec:milestone} and \S\ref{sec:res-learned} both find naive Monte Carlo winning on
	raw FOM at $L=16$, and argue this should be read as a finite-size effect rather than a
	failure of variance reduction, since naive's cost is independent of rarity while WE's
	penetration overhead is only worth paying once naive starts to struggle. We tested this
	directly with a lattice-size scan at fixed $T=2.6$, $\mathrm{beyond}=1$, $\tau=2$: $L\in
	\{8,16,24,32\}$, naive\_chains/naive\_sweeps held fixed across all four so that naive's
	cost is genuinely $L$-independent as the argument requires (Table~\ref{tab:lscaling}).
	
	% --- TABLE 8 ---
	\begin{table}[H]
		\centering\small
		\caption{EA spin glass, $\mathrm{beyond}=1$, naive vs.\ the best-performing WE variant at
			each $L$ (which variant wins changes with $L$; see text). $\hat\pi$ for naive shown to
			indicate how rare the target is at each scale.}
		\label{tab:lscaling}
		\begin{tabular}{lrrrr}
			\toprule
			$L$ & naive $\hat\pi$ & naive FOM & best WE FOM (method) & ratio \\
			\midrule
			8 & $1.79\times10^{-4}$ & $2.00\times10^{-6}$ & $1.04\times10^{-6}$ ($\mathrm{WE}[E]$) & $1.9\times$ \\
			16 & $5.18\times10^{-5}$ & $1.72\times10^{-7}$ & $4.83\times10^{-8}$ ($\mathrm{WE}[m]$) & $3.5\times$ \\
			24 & $1.30\times10^{-5}$ & $2.79\times10^{-8}$ & $8.06\times10^{-9}$ ($\mathrm{WE}[E]$) & $3.5\times$ \\
			32 & $4.17\times10^{-6}$ & $6.11\times10^{-9}$ & $1.69\times10^{-9}$ ($\mathrm{WE}[I_\theta]$) & $3.6\times$ \\
			\bottomrule
		\end{tabular}
	\end{table}
	
	Naive wins at every tested $L$, and the margin does not shrink; if anything it
	grows slightly and then plateaus, from $1.9\times$ at $L=8$ to $\sim\!3.5\times$ by
	$L=16$ and beyond (Figure~\ref{fig:lscaling}). This holds across a factor of ${\sim}40$ in target rarity
	($\hat\pi$ from $1.8\times10^{-4}$ down to $4.2\times10^{-6}$) and is robust to which WE
	variant is best at a given $L$ (the winner changes: $\mathrm{WE}[E]$ at $L=8,24$,
	$\mathrm{WE}[m]$ at $L=16$, $\mathrm{WE}[I_\theta]$ at $L=32$). This directly answers the
	open question of \S\ref{sec:milestone}: within the range tested, there is no sign
	of the naive/WE crossover opening as $L$ grows; if it exists at all on this problem, it is
	beyond $L=32$ at this temperature, not a nearby, easily-reached regime.
	
	% --- FIGURE 6 ---
	\noindent
	\begin{minipage}{\linewidth}
		\centering
		\includegraphics[width=0.92\textwidth]{figs/fig_lscaling.pdf}
		\captionof{figure}{Table~\ref{tab:lscaling} as a figure. Left: naive vs.\ best WE-variant FOM by
			$L$, log scale, labeled with which variant won. Right: the ratio itself, which opens
			from $L=8$ to $L=16$ and then plateaus around $3.5\times$, with no sign of narrowing
			back toward the naive/WE crossover within the tested range.}
		\label{fig:lscaling}
	\end{minipage}
	\par\vspace{0.8\baselineskip}
	
	One secondary pattern is visible but not claimed here beyond a point estimate:
	$\mathrm{WE}[I_\theta]$ is the weakest WE variant at $L=8,16,24$ but the \emph{strongest}
	at $L=32$, consistent with a learned coordinate's advantage over hand-picked ones
	mattering more as the landscape gets harder. We do not treat this as established: this
	scan predates the seeding fix of \S\ref{sec:milestone}, so each $\mathrm{WE}[I_\theta]$
	entry in Table~\ref{tab:lscaling} is one random draw from an uncharacterized distribution
	over training outcomes, the same seed-variance problem characterized in full in
	Table~\ref{tab:milestone}.
	Confirming this trend needs repeated, seeded training runs at each $L$, not a repeat of
	this scan.
	
	A boundary marker at $L=40$. An attempt one step further was stopped before
	$\mathrm{WE}[I_\theta]$ completed, but the partial result is informative on its own: at
	this size, \emph{naive itself} fails in 10 of 12 replicas, with $\mathrm{WE}[m]$ (11/12)
	and $\mathrm{WE}[E]$ (10/12) failing similarly badly. Per the ladder-difficulty discipline
	of \S\ref{sec:methods}, this is not a trustworthy comparison as run: every method is in
	the noise-over-noise regime simultaneously, the same failure mode that motivated that
	discipline in the first place, so we do not report it as a fifth point in
	Table~\ref{tab:lscaling}. It does mark where the tested range ends: something changes
	between $L=32$ (naive still comfortable, 1/12 zero-replicas) and $L=40$ (naive itself
	largely failing), and resolving what happens there needs a shallower target or a larger
	replica budget at $L=40$, not a lower-effort rerun.
	
	\subsection{FW-CADIS: real cost advantage, no flattening}\label{sec:fwcadis}
	The implementation is gated clean against exact enumeration at $L=4$ for both allocation
	settings, every threshold, both models: ratios of 0.999--0.998--0.989 (uniform) and
	1.005--1.012--0.992 (FW-CADIS, $\alpha=1$) on the ferromagnet, and 1.009--1.004--1.007 /
	0.986--0.980--0.983 on the spin glass, all within $\sim$2\% of exact for every threshold in
	the family, confirming that all thresholds are correctly tallied from a single run
	regardless of allocation, the core mechanic the whole method depends on.
	
	Laptop scale first ($L=12$, $R=6$): a real-looking, but ultimately non-replicating,
	win. Table~\ref{tab:fwall} collects every configuration tried across this section, from
	this first laptop-scale sweep through the final ferromagnet control; its top group sweeps
	$\alpha$ against uniform allocation on 6 thresholds spanning progress $0.668\to1.153$.
	
	% --- TABLE 9 ---
	\begin{table}[H]
		\centering\small
		\caption{FW-CADIS across every configuration tested, laptop scale through the ferromagnet
			control. Flatness marked $^*$ is computed only over thresholds actually observed, not
			comparable to unmarked rows in the same group. Dashes mark columns that do not apply to
			that group (coverage was 6/6 throughout past the laptop scale; a bootstrap CI was only
			computed for the final two configurations).}
		\label{tab:fwall}
		\begin{tabular}{lrrrr}
			\toprule
			Allocation & coverage & flatness (max/min) & 95\% CI & cost \\
			\midrule
			\multicolumn{5}{l}{\emph{Laptop scale}, $L=12$, $R=6$, 6 thresholds} \\
			uniform ($\alpha=0$) & 6/6 & 283.8 & --- & $2.01\times10^{7}$ \\
			$\alpha=0.25$ & 6/6 & 245.2 & --- & $7.87\times10^{6}$ \\
			$\alpha=0.5$ & 5/6 & $65.8^*$ & --- & $3.41\times10^{6}$ \\
			$\alpha=0.75$ & 5/6 & $63.5^*$ & --- & $3.29\times10^{6}$ \\
			$\alpha=1.0$ & 5/6 & $84.8^*$ & --- & $3.24\times10^{6}$ \\
			naive (reference) & 6/6 & 463.5 & --- & $1.61\times10^{7}$ \\
			\multicolumn{5}{l}{\emph{Full scale}, $L=16$, $R=30$, 6 thresholds (every allocation reaches 6/6 coverage)} \\
			uniform ($\alpha=0$) & 6/6 & 298.0 & --- & $1.55\times10^{8}$ \\
			$\alpha=0.25$ & 6/6 & 377.8 & --- & $7.89\times10^{7}$ \\
			$\alpha=0.5$ & 6/6 & 323.9 & --- & $2.35\times10^{7}$ \\
			$\alpha=0.75$ & 6/6 & 342.4 & --- & $1.98\times10^{7}$ \\
			$\alpha=1.0$ & 6/6 & 385.5 & --- & $1.94\times10^{7}$ ($\sim\!8\times$ cheaper) \\
			naive (reference) & 6/6 & 445.1 & --- & $1.34\times10^{8}$ \\
			\multicolumn{5}{l}{\emph{Final EA spin-glass configuration}, $L=16$, $R=90$, shallowest calibrated family} \\
			Uniform ($\alpha=0$) & --- & 276.1 & $[224,\,344]$ & $1.59\times10^{8}$ \\
			FW-CADIS ($\alpha=1$) & --- & 297.3 & $[234,\,388]$ & $1.99\times10^{7}$ ($\sim\!8\times$ cheaper) \\
			\multicolumn{5}{l}{\emph{Ferromagnet control}, $L=16$, $R=30$, shallowest calibrated family} \\
			Uniform ($\alpha=0$) & --- & 100.9 & $[65.8,\,160.3]$ & $1.96\times10^{8}$ \\
			FW-CADIS ($\alpha=1$) & --- & 287.3 & $[194.7,\,414.8]$ & $3.09\times10^{7}$ \\
			\bottomrule
		\end{tabular}
	\end{table}
	
	The only setting with full coverage and better flatness than uniform is $\alpha=0.25$: a
	real but modest $\sim$14\% improvement ($283.8\to245.2$) at $\sim$2.6$\times$ lower cost.
	Larger $\alpha$ looks dramatically flatter but is disqualified by the coverage-first rule
	of \S\ref{sec:methods}: it stops observing the deepest threshold at all, which would reward
	outright failure if scored on the surviving thresholds alone.
	
	Full scale ($L=16$, $R=30$) reverses the laptop-scale finding. The identical sweep
	at full scale (Table~\ref{tab:fwall}, second group) is the first sign that the method's
	behavior is scale-dependent in a way not anticipated from the smaller test.
	
	At full scale, \emph{every} tested $\alpha$ is worse than uniform on flatness, not better:
	the opposite of the laptop-scale result, and not even monotonic in $\alpha$ among the
	FW-CADIS runs ($\alpha=0.5$ beats both $0.25$ and $1.0$). All are substantially cheaper
	(2--8$\times$), so the method is not doing nothing, but it is not delivering the flattening
	it exists to demonstrate at this lattice size, bin count, and threshold spacing. This
	finding is what the rest of \S\ref{sec:fwcadis} exhaustively follows up on.
	
	Exhausting every remaining lever. Every lever available to improve on the full-scale
	finding above was tried and is reported in the third group of Table~\ref{tab:fwall}:
	bootstrap confidence
	intervals added to the flatness metric (not present at the smaller replica counts above); a
	replica increase from $R=30$ to $R=90$; wiring the learned coordinate of
	\S\ref{sec:milestone} in as the FW-CADIS binning axis instead of the physical progress
	coordinate (gated clean at $L=4$, tested at scale: no improvement, roughly double the
	cost); a common-random-numbers paired-difference design, which we confirmed does
	\emph{not} engage on this problem (correlation between the two allocations'
	same-seed estimates measured $\approx-0.05$ to $-0.10$, essentially zero, because the two
	allocations consume different numbers of walkers per bin and desynchronise the random
	stream almost immediately after the first resampling step); and laddering the threshold
	family to its shallowest calibrated depth, combined with the larger replica count.
	
	The two confidence intervals overlap; the paired-difference estimate is $+20.9$ with 95\%
	CI $[-77.2,\,+137.6]$, also including zero. Across every $\alpha$, replica count, and
	coordinate tried this session, the point-estimate \emph{direction} consistently favored
	uniform allocation over FW-CADIS on raw flatness, never the reverse. The observed
	confidence-interval width scaled almost exactly as $1/\sqrt{R}$ between $R=30$ and $R=90$
	(predicted shrink factor 0.577, observed 0.597); extrapolating that scaling, fully resolving
	the remaining gap would need on the order of $R\approx2300$ replicas ($\sim$8 hours of
	additional compute), and given the consistent point-estimate direction, such a run would
	most likely \emph{confirm} uniform allocation as measurably flatter rather than produce a
	win for FW-CADIS. We did not run it. The result we report instead: on this problem,
	FW-CADIS does not measurably flatten relative error across a threshold family better than
	uniform allocation, but it does so at approximately one-eighth the simulation cost for an
	otherwise statistically indistinguishable error profile. The cost advantage is real,
	reproducible, and does not require more replicas to state; the flattening advantage the
	method is named for is the part that did not materialize here.
	
	Is this specific to the adversarial spin-glass landscape? No, it is worse, and
	statistically confirmed, on the ``easy'' ferromagnet. Every result above uses the EA spin
	glass, the harder of the two models by design. To test whether the flattening failure is a
	property of that specific rugged landscape rather than a more general limit of the method,
	we ran the identical comparison on the ferromagnet, where the reaction coordinate is known
	to be good and the equilibrium landscape is smooth (Table~\ref{tab:fwall}, fourth group).
	
	Here the two confidence intervals do \emph{not} overlap: the paired-difference estimate is
	$+178.2$ with 95\% CI $[+59.8,\,+351.3]$, which excludes zero. On the easier,
	non-adversarial landscape, FW-CADIS is not merely inconclusive but measurably, statistically
	confirmed \emph{worse} than uniform allocation, the opposite of what a
	landscape-difficulty explanation would predict, since the spin glass (where the method's
	own hypothesis says a global map should matter most) never even reached this level of
	statistical resolution in either direction. Common-random-numbers correlation here measured
	$0.062$, consistent with zero and with the mechanism identified in \S\ref{sec:kimcai}: the
	two allocations consume different walker populations per bin regardless of which model is
	sampled. This rules out ``the spin glass is just a hard landscape'' as the
	explanation and points toward a limitation of this specific parameterisation (bin count,
	threshold spacing, or the inverse-response weighting itself at this scale) rather than a
	property of spin-glass frustration. Distinguishing which of those remains open, but the
	question of \emph{whether} the effect is landscape-specific is now answered: it is not.
	Figure~\ref{fig:fwcadis} plots the per-threshold error underlying both comparisons.
	
	% --- FIGURE 7 ---
	\noindent
	\begin{minipage}{\linewidth}
		\centering
		\includegraphics[width=0.92\textwidth]{figs/fig_fwcadis.pdf}
		\captionof{figure}{Per-threshold relative error (log scale), uniform vs.\ FW-CADIS allocation, at
			the final tested configuration for each model (Table~\ref{tab:fwall}, third group spin
			glass, fourth group ferromagnet). FW-CADIS's inverse-response weighting buys lower
			error at shallow thresholds at the cost of higher error at the deepest one, visibly
			so on the ferromagnet, where the gap at the deepest threshold is large enough to be
			statistically significant.}
		\label{fig:fwcadis}
	\end{minipage}
	\par\vspace{0.8\baselineskip}
	
	\subsection{What is established}
	
	The adjoint-committor identity is exact. Tilting by the exact committor
	is provably zero-variance; this was confirmed numerically to $\sim\!10^{-30}$ relative
	variance on a solvable toy walk before any lattice code was written (\S\ref{sec:identity}).
	
	Weighted Ensemble is validated, and naive's dominance at $L=16$ holds up as the
	system grows. Weighted Ensemble, weight windows applied to spin configurations,
	reproduces exact enumeration at $L=4$ to 1--2\%, and at $L=20$ reaches $\sim$2 decades
	deeper into a magnetization tail than naive Monte Carlo at matched cost
	(\S\ref{sec:res-we}). Naive Monte Carlo's dominance over every tested weight-window
	variant, not established at $L=16$ alone, held up across a four-point lattice-size scan
	up to $L=32$ with no sign of narrowing (\S\ref{sec:lscaling}), the one claim in this
	paper that grew \emph{more}, not less, confident under additional scrutiny.
	
	At a depth where naive is provably blind, the self-consistency-trained
	coordinate wins significantly. It finds the rare event
	significantly more often than the hand-picked energy coordinate: pooled across three
	independent networks, $66.7\%$ versus $40.0\%$ (Fisher exact $p=0.017$,
	\S\ref{sec:deeprare}, Figure~\ref{fig:deeprare}), the paper's one statistically
	significant positive result, and
	its clearest evidence that the transport toolkit's distinctive machinery earns its keep
	specifically where naive cannot compete.
	
	\subsection{What remains open}
	Naive dominance is depth-dependent, not universal, and the distinction matters for
	how every other result in this section should be read. It is robust to lattice size at a
	fixed, shallow threshold depth (\S\ref{sec:lscaling}: $L=8$--$32$, no sign of the margin
	closing), but not to depth itself at fixed size: pushed deep enough that naive fails
	outright (\S\ref{sec:deeprare}), the transport toolkit's advantage becomes real, and for the
	self-consistency coordinate, significant. The two scans probe different axes of the same
	underlying claim and give different answers: treating ``naive wins at $L=16$'' as a fixed,
	depth-independent property of the problem would have missed the paper's one significant
	positive result entirely.
	
	The learned coordinate trained on milestoned committor labels is the clearest case
	of an apparent win dissolving under scrutiny, and its sign is not resolved by seed count
	alone. A first training run beat the hand-picked
	energy coordinate by $+14\%$; seventeen retraining runs of the identical configuration gave
	outcomes from $-52\%$ to $+112\%$ (Table~\ref{tab:milestone}), and pooling them with a
	two-level bootstrap left a positive but non-significant point estimate, $1.19\times$ with
	95\% CI $[0.84,1.75]$ (\S\ref{sec:milestone}). Bin-range calibration and bin-count tuning
	were both tested and ruled out as the driver of that spread, which points to the variance
	being intrinsic to what each trained network learns rather than to how its output is
	binned. Closing the interval by seed count alone would need on the order of 100 seeds; a
	variance-reduction step applied to the training stage itself, rather than more independent
	draws from it, is the more economical route and remains the natural next test.
	
	FW-CADIS's flattening claim is not established, and the negative result is now
	well characterized rather than merely observed. Its cost advantage over uniform
	allocation is real and reproducible ($\sim\!8\times$ cheaper for a statistically
	indistinguishable error profile), but the flattening property it is named for did not
	appear at any $\alpha$, replica count, or coordinate tried, and a ferromagnet control
	(\S\ref{sec:fwcadis}) rules out spin-glass ruggedness as the explanation: if anything the
	failure is \emph{worse} on the easier landscape. What remains open is which of two
	explanations is responsible, the specific bin/threshold parameterisation or a more general
	limit of inverse-response weighting at this scale; distinguishing them needs a parameter
	sweep independent of $\alpha$, not another problem instance or more replicas.
	
	The self-consistency objective's FOM advantage over this project's own surrogate is
	not established at shallow depth (\S\ref{sec:kimcai}), and did not visibly grow with
	depth within the range where naive still competes, though its reliability advantage
	(fewest zero-replicas of any weight-window method, at every shallow depth tested) held
	throughout, foreshadowing the significant reliability result the same objective produced
	once tested where naive genuinely fails (\S\ref{sec:deeprare}).
	
	Push-to-scale questions (three dimensions, larger $L$, deeper than beyond=4)
	are entirely untouched, and are the natural next test of whether the deep-rare
	advantage found here continues to grow or plateaus.
	
	\subsection{Relation to prior art}\label{sec:pos}
	Kim \& Cai (arXiv:2602.12294, 2026) learn an importance function $I(s)$, modify
	transitions as $K'(i\to j)=K(i\to j)I(j)/I(i)$, and use a branching random walk: in the
	language of \S\ref{sec:identity}, their $I$ is our $h$, their $K'$ is the tilt, and their
	branching is splitting: they independently reconstruct the zero-variance/weight-window
	recipe for discrete rare transitions, without the transport framing, and without comparing
	their learned coordinate against a hand-picked reaction coordinate the way this paper does.
	We directly implemented and tested their self-consistency training objective
	(\S\ref{sec:kimcai}); what remains open after doing so is not whether it works in
	principle (their paper demonstrates convergence to the exact transition rate on two test
	potentials) but whether it offers a measurable practical advantage over simpler
	alternatives on a problem of this kind at this scale, which our data does not yet support.
	Doob-transform VAN (Phys.\ Rev.\ E 111, 034120, 2025) approximates the
	leading eigenvector of the tilted generator for \emph{dynamical} large deviations, the
	dynamical analogue of $\psi\dg$; a head-to-head against a reverse-KL VAN baseline, which
	this comparison would require, remains untried in this project. No successor work citing
	either paper closes the specific gap this paper's methods and results address, as of the
	most recent check.
	
	\FloatBarrier
	
	% ========================================================================
	\section{Conclusions}\label{sec:conclusions}
	The transport--statistical-mechanics correspondence this project set out to test is real
	where it can be proved (the adjoint$=$committor identity), but none of the more
	experimental claims built on top of it held up as cleanly as an initial pass through the
	results suggested. The clearest lesson is not any single number: ``0/10 zero-replicas''
	and a tidy point estimate are not, by themselves, evidence of a real effect. FW-CADIS's
	flattening claim, the Kim \& Cai comparison, and the milestoning result each needed a
	bootstrap confidence interval to reveal an apparent win as noise, and in the milestoning
	case a seventeen-seed sweep on top of it, since a single training run is not a
	reproducible measurement of a learned coordinate at all
	(\S\ref{sec:milestone}). Reaching these conclusions required
	diagnosing and correcting real implementation issues along the way: a resampling
	interval that silently degenerated a weight-window method to naive sampling, a
	mis-calibrated bin range, an under-determined training objective, and an unseeded
	network initialization that let a single run be mistaken for a stable result, each
	caught by an independent cross-check rather than a number simply looking wrong.
	
	One result stands apart from that pattern, and it sharpens the paper's central
	claim. The self-consistency objective, whose shallow-depth
	advantage had itself dissolved under scrutiny and whose training cost is a real practical
	tax, roughly $170\times$ that of the surrogate objective (\S\ref{sec:kimcai}), produced this paper's
	one statistically significant positive result once tested at a depth naive genuinely
	cannot reach (\S\ref{sec:deeprare}): the one apparent win in this project that grew
	\emph{stronger}, not weaker, under the same scrutiny that dissolved every other one. Read
	together, the transport toolkit's value is not diffuse across every regime this
	correspondence touches. It is concentrated specifically in the deep-rare regime radiation
	transport was built for, and requires a training objective shaped like the actual
	committor condition rather than a shallow, energy-adjacent proxy of it. We think this is
	the right way to report a project whose most interesting finding, more than once, was
	that an apparent win did not survive being checked properly, and whose one win that did
	survive only appeared once the problem was made hard enough to need it.
	
	% ========================================================================
	\section*{Acknowledgements}
	The author gratefully acknowledges the support provided by Fahad Bin Sultan University, Tabuk, Saudi Arabia, for providing a supportive academic environment and access to the computational facilities necessary to conduct this research. The institutional support and research infrastructure were essential for implementing the Monte Carlo experiments and performing the systematic statistical analysis reported in this work.
	
	% To print the CRediT authorship contribution details
	\printcredits
	
	% ========================================================================
	\appendix
	\renewcommand{\thetable}{A.\arabic{table}}
	\setcounter{table}{0}
	\section{Supporting tables}\label{app:tables}
	The point estimates collected here support the head-to-head comparison discussed in
	\S\ref{sec:kimcai}. They are reported in full for reproducibility; the corresponding
	confidence intervals, and the conclusions drawn from them, are given in the main text
	(Table~\ref{tab:kimcaiboot}, Figure~\ref{fig:kimcai}).
	
	% --- Kim & Cai point estimates ---
	\begin{table}[H]
		\centering\small
		\caption{EA spin glass, $L=16$, shallowest calibrated target, 10 replicas. Point estimates
			only; the bootstrapped confidence intervals are given in
			Table~\ref{tab:kimcaiboot} of the main text.}
		\label{tab:kimcai}
		\begin{tabular}{lrrrr}
			\toprule
			Method & $\hat\pi$ & rel.\ s.d.\ & FOM & zero-replicas \\
			\midrule
			naive & $3.36\times10^{-5}$ & 0.294 & $8.60\times10^{-8}$ & 0/10 \\
			$\mathrm{WE}[m]$ & $3.94\times10^{-5}$ & 0.553 & $8.81\times10^{-8}$ & 1/10 \\
			$\mathrm{WE}[E]$ & $3.39\times10^{-5}$ & 0.671 & $1.33\times10^{-8}$ & 0/10 \\
			$\mathrm{WE}[I_\theta]$ surrogate & $2.92\times10^{-5}$ & 0.904 & $8.62\times10^{-9}$ & 1/10 \\
			$\mathrm{WE}[I_\theta]$ self-consistent & $4.01\times10^{-5}$ & 0.655 &
			$1.62\times10^{-8}$ & 0/10 \\
			\bottomrule
		\end{tabular}
	\end{table}
	
	%% Bibliography
	\bibliographystyle{unsrtnat}
	\bibliography{cas-refs}
	
\end{document}